\documentclass[12pt,a4paper]{article}
\usepackage[utf8]{inputenc}
\usepackage[T1]{fontenc}
\usepackage[spanish]{babel}
\usepackage{geometry}
\usepackage{xcolor}
\usepackage{listings}
\usepackage{booktabs}
\usepackage{longtable}
\usepackage{array}
\usepackage{tikz}
\usetikzlibrary{shapes.geometric, arrows.meta, positioning, calc, fit, backgrounds}
\usepackage{hyperref}
\usepackage{fancyhdr}
\usepackage{titlesec}
\usepackage{amsmath}
\usepackage{enumitem}
\usepackage{tcolorbox}
\tcbuselibrary{skins, breakable}
\usepackage{alltt}

\geometry{top=2.5cm, bottom=2.5cm, left=2.8cm, right=2.8cm}
\setlength{\headheight}{15pt}

% -- Colors ------------------------------------------------------------------
\definecolor{primary}{RGB}{108,99,255}
\definecolor{secondary}{RGB}{167,139,250}
\definecolor{success}{RGB}{76,175,138}
\definecolor{codebg}{RGB}{248,249,250}
\definecolor{codegreen}{RGB}{0,128,0}
\definecolor{codered}{RGB}{197,15,31}

% -- Header/Footer -----------------------------------------------------------
\pagestyle{fancy}
\fancyhf{}
\fancyhead[L]{\textcolor{primary}{\textbf{CogniFace}} --- Documentacion Tecnica}
\fancyhead[R]{\textcolor{gray}{\small Universidad Tecnica del Norte}}
\fancyfoot[C]{\thepage}
\renewcommand{\headrulewidth}{0.4pt}

% -- Titles ------------------------------------------------------------------
\titleformat{\section}[block]{\large\bfseries\color{primary}}{\thesection}{1em}{}[\titlerule]
\titleformat{\subsection}[block]{\normalsize\bfseries\color{primary!80!black}}{\thesubsection}{1em}{}

% -- Code listings -----------------------------------------------------------
\lstdefinelanguage{JSCustom}{
  keywords={const,let,var,export,import,async,await,function,return,from,if,else,
            case,switch,true,false,null,undefined,new,class,this},
  sensitive=true,
  comment=[l]{//},
  morecomment=[s]{/*}{*/},
  morestring=[b]',
  morestring=[b]"
}
\lstdefinestyle{jsStyle}{
  backgroundcolor=\color{codebg},
  commentstyle=\color{codegreen},
  keywordstyle=\color{primary}\bfseries,
  numberstyle=\tiny\color{gray},
  stringstyle=\color{codered},
  basicstyle=\ttfamily\small,
  breaklines=true,
  keepspaces=true,
  numbers=left,
  numbersep=5pt,
  showstringspaces=false,
  tabsize=2,
  frame=single,
  rulecolor=\color{primary!30!white},
  language=JSCustom
}
\lstset{style=jsStyle}

% -- Boxes -------------------------------------------------------------------
\tcbset{
  infobox/.style={enhanced, breakable,
    colback=primary!5!white, colframe=primary!50!white,
    fonttitle=\bfseries\color{primary}, arc=3mm},
  codebox/.style={enhanced, colback=codebg, colframe=gray!30,
    fonttitle=\bfseries\small, arc=2mm},
  successbox/.style={enhanced, breakable,
    colback=success!8!white, colframe=success!60!white,
    fonttitle=\bfseries\color{success!70!black}, arc=3mm}
}

% ============================================================================
\begin{document}

% -- Portada -----------------------------------------------------------------
\begin{titlepage}
\begin{tikzpicture}[remember picture, overlay]
  \fill[primary!10!white] (current page.north west)
    rectangle ([yshift=-5.5cm]current page.north east);
\end{tikzpicture}

\vspace*{2.5cm}
\begin{center}
  {\LARGE\bfseries\textcolor{primary}{UNIVERSIDAD TECNICA DEL NORTE}}\\[0.4cm]
  {\large Facultad de Educacion, Ciencia y Tecnologia (FECYT)}\\[0.2cm]
  {\large Carrera de Psicologia}\\[2cm]

  \begin{tcolorbox}[infobox, width=0.92\textwidth]
    \centering
    {\Huge\bfseries\textcolor{primary}{CogniFace}}\\[0.5cm]
    {\LARGE Documentacion Tecnica del Sistema}\\[0.3cm]
    {\large Plataforma Web para Experimento N-Back con Rostros}
  \end{tcolorbox}

  \vspace{2cm}
  \begin{tabular}{rl}
    \textbf{Tipo de documento:} & Informe Tecnico v1.0 \\[0.3cm]
    \textbf{Repositorio:}       & \texttt{zeuspyEC/CogniFace} \\[0.3cm]
    \textbf{URL publica:}       & \texttt{https://cogniface.web.app} \\[0.3cm]
    \textbf{Fecha:}             & Junio 2026 \\[0.3cm]
    \textbf{Stack:}             & React 18 + Firebase Spark
  \end{tabular}
\end{center}
\end{titlepage}

\tableofcontents
\newpage

% ============================================================================
\section{Resumen Ejecutivo}

CogniFace es una aplicacion web de una sola pagina (SPA) disenada para administrar
el experimento N-Back con rostros de la Carrera de Psicologia de la UTN-FECYT.
Permite recolectar datos de memoria de trabajo sobre participantes anonimos,
calcular el Indice de Afinidad Facial (IAF) y visualizar resultados desde
un panel administrativo protegido.

\begin{tcolorbox}[infobox, title={Caracteristicas principales}]
\begin{itemize}[leftmargin=1em]
  \item Experimento N-Back de una sola fase (N=1 o N=2) elegida por el investigador
  \item 5 ensayos de practica con retroalimentacion + 20 experimentales sin feedback
  \item Medicion de tiempo de reaccion sub-milisegundo con \texttt{performance.now()}
  \item Persistencia en Firestore (Spark gratuito) con escritura atomica al finalizar
  \item Panel admin con graficas estadisticas, tabla CRUD y exportacion a Excel
  \item Modo oscuro / claro con variables CSS y \texttt{localStorage}
  \item CI/CD automatico via GitHub Actions a Firebase Hosting
\end{itemize}
\end{tcolorbox}

% ============================================================================
\section{Stack Tecnologico}

\begin{longtable}{>{\bfseries}p{3.5cm} p{4.8cm} p{4.8cm}}
\toprule
Capa & Tecnologia & Proposito \\
\midrule
\endfirsthead
\toprule
Capa & Tecnologia & Proposito \\
\midrule
\endhead
\bottomrule
\endfoot
Frontend  & React 18 + Vite 6   & SPA con hot-reload y bundling optimizado \\
Enrutado  & React Router v6     & Rutas \texttt{/} (experimento) y \texttt{/admin} \\
Estilos   & CSS variables       & Dark/light mode via \texttt{:root} y \texttt{[data-theme]} \\
Graficas  & Recharts 2.x        & Slope chart, lollipop, comparacion de grupos \\
Excel     & SheetJS (xlsx)      & Exportacion en cliente, sin servidor \\
Auth      & Firebase Auth       & Login email/password para investigador \\
BD        & Firestore (Spark)   & Colecciones \texttt{participants} + \texttt{trials} \\
Hosting   & Firebase Hosting    & CDN global, SSL automatico \\
CI/CD     & GitHub Actions      & Build + deploy en cada push a \texttt{main} \\
\bottomrule
\end{longtable}

% ============================================================================
\section{Arquitectura del Sistema}

\subsection{Vista General}

\begin{center}
\begin{tikzpicture}[
  node distance=1.5cm,
  box/.style={rectangle, draw=primary, fill=primary!8, rounded corners=4pt,
              text width=2.8cm, align=center, minimum height=0.9cm, font=\small\bfseries},
  arr/.style={-{Stealth[length=5pt]}, thick, color=primary!70},
  lbl/.style={font=\scriptsize\color{gray}, midway}
]
\node[box] (browser) {Navegador\\(React SPA)};
\node[box, right=3.5cm of browser] (auth) {Firebase Auth};
\node[box, above=1.2cm of auth] (hosting) {Firebase Hosting};
\node[box, below=1.2cm of auth] (firestore) {Firestore DB};
\node[box, left=3.5cm of hosting] (ci) {GitHub Actions};

\draw[arr] (hosting) -- node[lbl, above]{CDN/HTTPS} (browser);
\draw[arr] (browser) -- node[lbl, above]{login} (auth);
\draw[arr] (browser) -- node[lbl, right]{read/write} (firestore);
\draw[arr] (ci)      -- node[lbl, above]{deploy} (hosting);

\begin{scope}[on background layer]
  \node[draw=secondary!40, fill=secondary!5, rounded corners=6pt,
        fit=(auth)(firestore)(hosting), inner sep=8pt] {};
\end{scope}
\end{tikzpicture}
\end{center}

\subsection{Estructura de Archivos}

\begin{tcolorbox}[codebox, title={Archivos principales del proyecto}]
\begin{alltt}
\small
src/
  components/
    experiment/
      WelcomeScreen.jsx      Portada UTN/FECYT
      ParticipantForm.jsx    Formulario sexo/ID/edad
      BlockSelector.jsx      Eleccion N1 o N2
      InstructionsScreen.jsx Instrucciones por bloque
      PracticeBlock.jsx      5 ensayos con feedback
      ExperimentBlock.jsx    20 ensayos sin feedback
      ResultsScreen.jsx      Metricas + IAF
    admin/
      AdminDashboard.jsx     Layout header + tabs
      AdminLogin.jsx         Login Firebase Auth
      ParticipantsTable.jsx  CRUD + Export Excel
      ChartsPanel.jsx        Graficas Recharts
    shared/
      ThemeToggle.jsx        Boton dark/light
      ProtectedRoute.jsx     Guard /admin
  hooks/
    useExperimentEngine.js   Maquina de estados
  context/
    ThemeContext.jsx
    AuthContext.jsx
    ExperimentContext.jsx
  lib/
    sequences.js             Generador N-Back
    statistics.js            IAF, metricas
    firestoreService.js      CRUD Firestore
    imagePreloader.js        Precarga 12 rostros
\end{alltt}
\end{tcolorbox}

% ============================================================================
\section{Maquina de Estados del Experimento}

El hook \texttt{useExperimentEngine.js} implementa una maquina de estados finita
con \texttt{useReducer}. Los timers se gestionan con \texttt{setTimeout} en efectos
React, garantizando que cada fase se active exactamente una vez por ensayo.

\begin{center}
\begin{tikzpicture}[
  state/.style={circle, draw=primary, fill=primary!12, thick,
                minimum size=1.8cm, font=\footnotesize\bfseries, align=center},
  arr/.style={-{Stealth[length=6pt]}, thick, color=primary!80},
  lbl/.style={font=\scriptsize\color{gray}, fill=white, inner sep=1pt}
]
\node[state] (idle)     at (0,0)      {idle};
\node[state] (load)     at (3,0)      {loading};
\node[state] (fix)      at (6,0)      {fixation};
\node[state] (stim)     at (6,-2.8)   {stimulus};
\node[state] (resp)     at (3,-2.8)   {response};
\node[state] (done)     at (0,-2.8)   {done};

\draw[arr] (idle) -- node[lbl, above]{start()} (load);
\draw[arr] (load) -- node[lbl, above]{LOADED} (fix);
\draw[arr] (fix)  -- node[lbl, right]{500 ms} (stim);
\draw[arr] (stim) -- node[lbl, below]{1000 ms} (resp);
\draw[arr] (resp) to[bend left=30] node[lbl, left]{END TRIAL} (fix);
\draw[arr] (resp) -- node[lbl, above]{ultimo} (done);
\end{tikzpicture}
\end{center}

\begin{tcolorbox}[infobox, title={Temporizado por ensayo (3500 ms totales)}]
\centering
\begin{tabular}{clc}
\toprule
Fase & Descripcion & Duracion \\
\midrule
Cruz (+) & Punto de fijacion & 500 ms \\
Rostro   & Estimulo facial   & 1\,000 ms \\
Ventana  & Respuesta ESPACIO & 2\,000 ms \\
\bottomrule
\end{tabular}
\end{tcolorbox}

\subsection{Reducer (fragmento clave)}

\begin{lstlisting}[caption={Caso RESPOND en useExperimentEngine.js}]
case 'RESPOND': {
  if (state.respondedThisTrial || state.phase !== 'response') return state
  const rt = performance.now() - state.stimulusOnset
  const trial = state.sequence[state.trialIndex]
  const immediateResult = trial?.is_target ? 'correct' : 'incorrect'
  return {
    ...state,
    respondedThisTrial: true,
    pendingRT: rt,
    lastFeedback: immediateResult  // visible en PracticeBlock
  }
}
\end{lstlisting}

\begin{tcolorbox}[infobox, title={Por que useReducer y no useState?}]
\texttt{useReducer} garantiza transiciones atomicas: el estado del ensayo
(fase, indice, RT, feedback) actualiza en una sola operacion de despacho,
eliminando race conditions entre los timers de fijacion, estimulo y respuesta.
\end{tcolorbox}

% ============================================================================
\section{Flujo de Datos del Participante}

\begin{center}
\begin{tikzpicture}[
  flowbox/.style={rectangle, draw=primary, fill=primary!8, rounded corners=3pt,
               minimum width=2.5cm, minimum height=0.8cm, font=\small\bfseries},
  db/.style={rectangle, draw=secondary, fill=secondary!10, rounded corners=3pt,
             minimum width=2.5cm, minimum height=0.8cm, font=\footnotesize},
  arr/.style={-{Stealth[length=5pt]}, thick, color=primary!70},
  darr/.style={-{Stealth[length=5pt]}, dashed, color=secondary}
]
\node[flowbox] (w) at (0,0)    {Welcome};
\node[flowbox] (f) at (0,-1.2) {Formulario};
\node[flowbox] (b) at (0,-2.4) {Block N1/N2};
\node[flowbox] (i) at (0,-3.6) {Instrucciones};
\node[flowbox] (p) at (0,-4.8) {Practica x5};
\node[flowbox] (e) at (0,-6.0) {Experimento x20};
\node[flowbox] (r) at (0,-7.2) {Resultados};

\node[db] (fs1) at (4,-1.2) {createParticipant};
\node[db] (fs2) at (4,-6.0) {batch write trials};

\draw[arr] (w)--(f); \draw[arr] (f)--(b); \draw[arr] (b)--(i);
\draw[arr] (i)--(p); \draw[arr] (p)--(e); \draw[arr] (e)--(r);
\draw[darr] (f) -- (fs1);
\draw[darr] (e) -- (fs2);
\end{tikzpicture}
\end{center}

% ============================================================================
\section{Esquema de Base de Datos Firestore}

\begin{longtable}{>{\ttfamily\small}p{4.5cm} >{\small}p{3cm} >{\small}p{5.5cm}}
\toprule
\normalfont\textbf{Campo} & \normalfont\textbf{Tipo} & \normalfont\textbf{Descripcion} \\
\midrule
\endfirsthead
\multicolumn{3}{l}{\small\textit{Coleccion: /participants/\{id\}}}\\
\toprule
\normalfont\textbf{Campo} & \normalfont\textbf{Tipo} & \normalfont\textbf{Descripcion} \\
\midrule
\endhead
\multicolumn{3}{l}{\small\textit{Continua...}}\\
\bottomrule
\endfoot
\bottomrule
\multicolumn{3}{l}{}
\endlastfoot
\multicolumn{3}{l}{\textbf{Coleccion: /participants/\{id\}}}\\[0.3em]
gender           & string  & "male" | "female" \\
participant\_code & string & ID asignado por investigador \\
age              & number  & Edad del participante \\
timestamp        & Timestamp & Hora de inicio \\
completed        & boolean & Experimento finalizado \\
n\_back          & number  & 1 o 2 \\
iaf              & number  & Indice de Afinidad Facial \\
iaf\_n1 / iaf\_n2 & number & IAF por condicion \\
hits / misses    & number  & Aciertos y omisiones \\
false\_alarms    & number  & Falsas alarmas \\
emotional\_errors & number & Errores en sexo opuesto \\
mean\_rt         & number  & TR promedio en ms \\
accuracy         & number  & Precision [0,1] \\[0.5em]
\multicolumn{3}{l}{\textbf{Subcoleccion: /participants/\{id\}/trials/\{id\}}}\\[0.3em]
block            & number  & 1 o 2 \\
trial\_number    & number  & Posicion en la secuencia \\
is\_practice     & boolean & Ensayo de practica \\
face\_id         & string  & "f01"..."m06" \\
face\_gender     & string  & "female" | "male" \\
is\_target       & boolean & Es un target N-Back \\
responded        & boolean & Participante presiono ESPACIO \\
reaction\_time   & number  & RT en ms (null si no respondio) \\
accuracy         & number  & 0 o 1 \\
error\_type      & string  & hit, miss, false\_alarm, correct\_rejection \\
\end{longtable}

% ============================================================================
\section{Sistema de Temas Dark/Light}

\begin{tcolorbox}[codebox, title={src/index.css --- Variables CSS}]
\begin{verbatim}
:root {
  --color-bg:         #0D1A2B;
  --color-surface:    #0F2035;
  --color-surface2:   #152D4A;
  --color-text:       #E8EFF9;
  --color-text-muted: #8892A4;
}

[data-theme="light"] {
  --color-bg:         #F0F4F8;
  --color-surface:    #FFFFFF;
  --color-surface2:   #E8EEF5;
  --color-text:       #1A2D42;
  --color-text-muted: #64748B;
}
\end{verbatim}
\end{tcolorbox}

El atributo \texttt{data-theme} se escribe en \texttt{document.documentElement}
desde \texttt{ThemeContext} y se persiste en \texttt{localStorage}.
El boton \texttt{ThemeToggle} aparece en las 7 interfaces de la app.

% ============================================================================
\section{Formula IAF}

\begin{tcolorbox}[infobox, title={Indice de Afinidad Facial (IAF)}]
\begin{align}
  \text{Hombre:}\quad & \text{IAF} =
    \text{acc}(F) - \text{acc}(M) \\[4pt]
  \text{Mujer:}\quad & \text{IAF} =
    \text{acc}(M) - \text{acc}(F)
\end{align}
donde $F$ = rostros femeninos, $M$ = rostros masculinos (solo ensayos reales).

\medskip
$\text{IAF} > 0 \;\Rightarrow\;$ ventaja hacia sexo opuesto $\;\Rightarrow\;$
hipotesis \textbf{apoyada}
\end{tcolorbox}

\begin{lstlisting}[caption={calculateIAF en statistics.js}]
export function calculateIAF(trials, participantGender) {
  const real = trials.filter(t => !t.is_practice)
  const accFemale = calculateAccuracy(real, 'female')
  const accMale   = calculateAccuracy(real, 'male')
  return participantGender === 'male'
    ? accFemale - accMale
    : accMale - accFemale
}
\end{lstlisting}

% ============================================================================
\section{Panel Administrativo}

\begin{longtable}{>{\bfseries}p{4cm} p{9cm}}
\toprule
Componente & Funcionalidad \\
\midrule
\endfirsthead
\toprule
Componente & Funcionalidad \\
\midrule
\endhead
\bottomrule
\endfoot
AdminLogin.jsx       & Login con Firebase Auth \\
AdminDashboard.jsx   & Header sticky con tabs, responsive (useMobile) \\
ParticipantsTable.jsx & Tabla paginada, filtros, boton Exportar Excel \\
ChartsPanel.jsx      & Dashboard de estadisticas (Recharts) \\
IAFWidget            & Comparacion IAF Hombres vs. Mujeres por bloque \\
SlopeChart           & Evolucion IAF por participante \\
GroupedBar           & Comparacion precision mismo/opuesto sexo \\
\end{longtable}

\subsection{Exportacion Excel (SheetJS)}

\begin{lstlisting}[caption={Funcion exportToExcel en ParticipantsTable.jsx}]
import * as XLSX from 'xlsx'

function exportToExcel(participants) {
  const wb = XLSX.utils.book_new()
  const rows = participants.map(p => ({
    'ID Participante': p.participant_code ?? '-',
    'Sexo':    p.gender === 'male' ? 'Masculino' : 'Femenino',
    'Edad':    p.age ?? '-',
    'N-Back':  p.n_back ?? '-',
    'IAF':     p.iaf?.toFixed(4) ?? '-',
    'Precision': p.accuracy ? (p.accuracy*100).toFixed(1)+'%' : '-',
  }))
  XLSX.utils.book_append_sheet(wb,
    XLSX.utils.json_to_sheet(rows), 'Participantes')
  XLSX.writeFile(wb, 'cogniface_participantes.xlsx')
}
\end{lstlisting}

% ============================================================================
\section{Pipeline CI/CD}

\begin{center}
\begin{tikzpicture}[
  flowbox/.style={rectangle, draw=primary, fill=primary!8, rounded corners=3pt,
               minimum width=2.5cm, minimum height=0.9cm,
               font=\small\bfseries, align=center},
  arr/.style={-{Stealth[length=5pt]}, thick, color=primary!70}
]
\node[flowbox] (push)   {git push main};
\node[flowbox, right=1.8cm of push]  (ci)     {GitHub Actions};
\node[flowbox, right=1.8cm of ci]    (build)  {npm run build};
\node[flowbox, right=1.8cm of build] (deploy) {firebase deploy};
\node[flowbox, below=1cm of deploy,  fill=success!12, draw=success]
  (live) {cogniface.web.app LIVE};

\draw[arr] (push)--(ci);
\draw[arr] (ci)--(build);
\draw[arr] (build)--(deploy);
\draw[arr] (deploy)--(live);
\end{tikzpicture}
\end{center}

El secreto \texttt{FIREBASE\_SERVICE\_ACCOUNT\_COGNIFACE} se configura en
GitHub Settings $>$ Secrets. El workflow esta en
\texttt{.github/workflows/deploy.yml}.

% ============================================================================
\section{Seguridad}

\begin{longtable}{>{\bfseries}p{4.2cm} p{9cm}}
\toprule
Aspecto & Implementacion \\
\midrule
\endfirsthead
\toprule
Aspecto & Implementacion \\
\midrule
\endhead
\bottomrule
\endfoot
Variables de entorno & Prefijo \texttt{VITE\_}, en \texttt{.env} (gitignoreado) \\
Firestore Rules      & Anonimos: solo crear docs en \texttt{participants}/\texttt{trials} \\
Auth admin           & Firebase Auth, solo investigador autenticado \\
ProtectedRoute       & Guard que redirige \texttt{/admin} si no autenticado \\
Sin Cloud Functions  & Toda logica en cliente, sin superficie de servidor \\
Sin PII              & Codigo anonimo, sexo, edad --- sin nombre \\
\end{longtable}

% ============================================================================
\section{Rendimiento y Compatibilidad}

\begin{itemize}
  \item \textbf{Imagenes como assets estaticos:} Las 12 fotografias se importan
        directamente en el bundle de Vite, eliminando latencia de red en runtime.
  \item \textbf{Bundle size:} 1\,367 kB (405 kB gzip). El tamano se debe a
        xlsx y Recharts; aceptable para uso experimental.
  \item \textbf{Timing:} \texttt{performance.now()} ofrece resolucion
        sub-milisegundo para medicion psicofisica de RT.
  \item \textbf{Soporte movil:} \texttt{ontouchstart}/\texttt{maxTouchPoints}
        cambia instrucciones ESPACIO a boton tactil; layout adaptable con
        \texttt{useMobile}.
\end{itemize}

% ============================================================================
\section*{Conclusion Tecnica}

CogniFace logra implementar un experimento cognitivo de precision milisegundo
con presupuesto de infraestructura cero (Firebase Spark gratuito), usando React
como maquina de estados experimental y Firestore como almacen longitudinal.
Build de produccion verificado: \texttt{0 errores, 667 modulos, 1.08 s}.

\vspace{2em}
\noindent\rule{\textwidth}{0.4pt}\\[0.3em]
{\small CogniFace v1.0 --- Universidad Tecnica del Norte,
FECYT, Carrera de Psicologia --- Junio 2026}

\end{document}
